\documentclass[times]{san}

\usepackage{array}

\bibliografia{bibliografia.bib}

% Pozwala łamać długie adresy URL w bibliografii także po cyfrach i literach.
\setcounter{biburlnumpenalty}{9000}
\setcounter{biburllcpenalty}{7000}
\setcounter{biburlucpenalty}{8000}

% Znacznik fragmentów, które zostaną uzupełnione w kolejnych etapach (PLAN.md, 13.2).
% Przed oddaniem: grep -r "douzupelnienia" main.tex rozdzialy/ musi zwracać pusty wynik.
\newcommand{\douzupelnienia}[1]{\textit{[Do uzupełnienia: #1]}}

\uczelnia{Społeczna Akademia Nauk}
\przedmiot{Systemy Szkieletowe}
\autor{Piotr Roman}
\nralbumu{122217}
\grupa{4}
\tytul{Aplikacja kierowcy -- system zarządzania dokumentami przewozowymi}
\prowadzacy{mgr inż. Krystian Gumiński}
\miasto{Łódź}
\rok{2026}

\begin{document}

\maketitlepage
\spistresci

\section{Wstęp}
\label{sec:wstep}

Kierowca w transporcie drogowym przez cały czas pracy korzysta z dokumentów
przewozowych. Najważniejszym z nich jest międzynarodowy list przewozowy CMR,
którego rolę i zawartość określa Konwencja o umowie międzynarodowego przewozu
drogowego towarów \cite{konwencja_cmr}. W przewozach krajowych i przy odbiorze
towaru z magazynu dochodzą do tego dokumenty wydania zewnętrznego (WZ),
faktury i inne potwierdzenia. Większość z nich nadal ma formę papierową:
kierowca otrzymuje je przy załadunku, zbiera pieczątki i podpisy przy rozładunku,
a potem dostarcza oryginały do biura firmy przewozowej.

Papierowy obieg dokumentów powoduje typowe problemy. Dokumenty giną lub niszczą
się w kabinie pojazdu, biuro poznaje ich treść dopiero po powrocie kierowcy,
a rozliczenie zlecenia i wystawienie faktury się opóźnia. Kierowca ma jednak
przy sobie telefon z aparatem, więc może dokument sfotografować i przesłać go
zaraz po otrzymaniu.

Niniejsza praca opisuje projekt i implementację \emph{Aplikacji kierowcy}:
responsywnej aplikacji webowej (PWA), w której kierowca fotografuje dokument
telefonem, kadruje i kompresuje zdjęcie, a następnie przechowuje je
w zaszyfrowanej postaci, z podziałem na typy, statusy i wersje.
Rozdział~\ref{sec:cel} określa cel projektu i wymagania, rozdział~\ref{sec:przeglad}
przedstawia istniejące rozwiązania, a rozdziały~\ref{sec:architektura}
i~\ref{sec:implementacja} opisują architekturę i implementację systemu.
Pracę zamyka podsumowanie z wnioskami (rozdział~\ref{sec:podsumowanie}).

\section{Cel projektu}
\label{sec:cel}

Celem projektu jest aplikacja, dzięki której kierowca ma dokumenty przewozowe
w telefonie, a nie tylko w papierowej teczce. Aplikacja ma skrócić czas od
otrzymania dokumentu do jego dostępności w formie cyfrowej, zmniejszyć ryzyko
zgubienia dokumentu i zapewnić jego poufność.

\subsection{Wymagania funkcjonalne}

\begin{enumerate}
  \item Rejestracja konta (formularz z walidacją danych) i logowanie loginem
        i hasłem; uwierzytelnianie żądań tokenem JWT.
  \item Profil kierowcy, do którego przypisane są jego dokumenty.
  \item Zarządzanie dokumentami: typ dokumentu (np. CMR, WZ), status
        (np. roboczy, przesłany, zaakceptowany), tytuł i numer, wyszukiwanie
        i filtrowanie.
  \item Zdjęcie dokumentu aparatem urządzenia mobilnego
        (interfejs \texttt{getUserMedia}), podgląd i kadrowanie przed zapisaniem
        oraz kompresja zdjęcia przed wysłaniem na serwer.
  \item Wersjonowanie dokumentów: każda nowa wersja pliku jest zachowywana,
        a zmiany metadanych i statusu trafiają do historii zmian.
\end{enumerate}

\subsection{Wymagania niefunkcjonalne}

\begin{enumerate}
  \item \textbf{Bezpieczeństwo}: pliki dokumentów są przechowywane
        w postaci zaszyfrowanej, hasła jako skróty kryptograficzne, a użytkownik
        ma dostęp wyłącznie do własnych dokumentów.
  \item \textbf{Obsługa na urządzeniach mobilnych}: interfejs responsywny,
        dostosowany do obsługi dotykowej, możliwy do zainstalowania jako PWA.
  \item \textbf{Łatwe wdrożenie}: cały system uruchamiany jednym poleceniem
        przez Docker Compose.
  \item \textbf{Niezawodność i obserwowalność}: monitorowanie stanu zdrowia
        usług, metryki, alerty i automatyczne restarty po awarii.
  \item \textbf{Jakość}: testy funkcjonalne każdej funkcji systemu,
        uzupełnione testami jednostkowymi.
\end{enumerate}

\section{Przegląd istniejących rozwiązań}
\label{sec:przeglad}

\douzupelnienia{analiza trzech rozwiązań według wspólnych kryteriów (funkcje,
platforma, obsługa e-CMR, skanowanie dokumentów, cena, wady) i tabela
porównawcza z projektowaną aplikacją; PLAN.md, punkt 13.4}

\subsection{Rozwiązanie A}
\douzupelnienia{np. platforma e-CMR}

\subsection{Rozwiązanie B}
\douzupelnienia{np. giełda transportowa z obiegiem dokumentów}

\subsection{Rozwiązanie C}
\douzupelnienia{np. aplikacja do skanowania dokumentów telefonem}

\section{Opis architektury}
\label{sec:architektura}

System składa się z aplikacji klienckiej działającej w przeglądarce (PWA),
serwera REST API, relacyjnej bazy danych oraz magazynu obiektów na pliki
dokumentów. Wszystkie komponenty działają w kontenerach Docker i są
uruchamiane wspólnie przez Docker Compose.

Rysunek~\ref{fig:kontenery} przedstawia kontenery systemu w notacji modelu C4.
Użytkownik łączy się wyłącznie z kontenerem \texttt{frontend}:
serwer nginx udostępnia zbudowaną aplikację React, a żądania na ścieżki
\texttt{/api/*} przekazuje do kontenera \texttt{backend}. Przeglądarka rozmawia
więc z jednym adresem (ta sama domena dla aplikacji i API), dzięki czemu nie
trzeba konfigurować mechanizmu CORS. Backend zapisuje dane w bazie PostgreSQL
i zaszyfrowane pliki w magazynie MinIO, zgodnym z interfejsem Amazon S3.
Prometheus okresowo pobiera metryki backendu, a Grafana prezentuje je na
dashboardach.

\begin{figure}[htbp]
\centering
\includegraphics[width=0.85\textwidth]{diagramy/kontenery.png}
\caption{Diagram kontenerów systemu (model C4)}
\label{fig:kontenery}
\zrodlo{opracowanie własne}
\end{figure}

Kontenery są rozdzielone na trzy sieci Dockera według zasady najmniejszych
uprawnień: \texttt{app} (nginx i backend), \texttt{data} (backend, PostgreSQL,
MinIO) oraz \texttt{monitoring} (backend, Prometheus, Grafana). Frontend nie ma
bezpośredniego dostępu do bazy danych ani do magazynu plików, a jedynym
kontenerem obecnym we wszystkich sieciach jest backend. Szczegóły konfiguracji
opisano w podrozdziale~\ref{sec:docker}.

\subsection{Wybór technologii}

W całym projekcie użyto języka TypeScript. Frontend i backend
korzystają dzięki temu z tych samych typów i schematów walidacji, umieszczonych
we wspólnym pakiecie. Zestawienie technologii zawiera
tabela~\ref{tab:technologie}.

\begin{table}[htbp]
\centering
\caption{Technologie użyte w projekcie}
\label{tab:technologie}
\small
\begin{tabular}{|p{0.2\textwidth}|p{0.3\textwidth}|p{0.38\textwidth}|}
\hline
\textbf{Obszar} & \textbf{Technologia} & \textbf{Uzasadnienie} \\
\hline
Backend & Node.js 22, Express & Wskazany w wymaganiach, prosty serwer REST \\
\hline
Baza danych & PostgreSQL 16 & Wskazana w wymaganiach \\
\hline
ORM i migracje & Prisma & Schemat bazy w jednym pliku, migracje wersjonowane w repozytorium \\
\hline
Pliki & MinIO (S3) & Magazyn obiektów w osobnym kontenerze \\
\hline
Walidacja & Zod & Te same reguły na froncie i w API \\
\hline
Frontend & React 19, Vite & Szybkie środowisko deweloperskie, łatwa konfiguracja PWA \\
\hline
Uwierzytelnianie & JWT, argon2 & Tokeny z wymagań; aktualny standard skracania haseł \\
\hline
Szyfrowanie & AES-256-GCM (\texttt{node:crypto}) & Szyfrowanie kopertowe plików przed zapisem \\
\hline
Monitoring & Prometheus, Grafana & Metryki, stan zdrowia, alerty, dashboardy \\
\hline
Konteneryzacja & Docker, Docker Compose & Uruchomienie całego systemu jednym poleceniem \\
\hline
Testy & Vitest, Supertest, Playwright & Testy funkcjonalne API i interfejsu, testy jednostkowe \\
\hline
\end{tabular}
\zrodlo{opracowanie własne}
\end{table}

\subsection{Struktura repozytorium}

Kod jest zorganizowany jako monorepo oparte na mechanizmie \emph{workspaces}
menedżera npm. Jedno polecenie \texttt{npm install} instaluje zależności
wszystkich pakietów, a wspólne skrypty (\texttt{build}, \texttt{lint},
\texttt{test}) uruchamiane są z katalogu głównego. Podział katalogów
przedstawia tabela~\ref{tab:repozytorium}.

\begin{table}[htbp]
\centering
\caption{Katalogi repozytorium}
\label{tab:repozytorium}
\small
\begin{tabular}{|p{0.22\textwidth}|p{0.66\textwidth}|}
\hline
\textbf{Katalog} & \textbf{Zawartość} \\
\hline
\texttt{backend/} & Serwer REST API (Express) \\
\hline
\texttt{frontend/} & Aplikacja kliencka React (PWA) \\
\hline
\texttt{shared/} & Schematy walidacji, typy i wyliczenia wspólne dla frontu i backendu \\
\hline
\texttt{e2e/} & Funkcjonalne testy interfejsu (Playwright) \\
\hline
\texttt{monitoring/} & Konfiguracja Prometheusa, Grafany i alertów \\
\hline
\texttt{docs/} & Niniejsze sprawozdanie (\LaTeX) i diagramy \\
\hline
\texttt{scripts/} & Skrypty pomocnicze, m.in. budowanie sprawozdania w Dockerze \\
\hline
\end{tabular}
\zrodlo{opracowanie własne}
\end{table}

Pakiet \texttt{shared} jest kompilowany jako pierwszy, ponieważ importują go
backend i frontend. Jakość kodu pilnują ESLint (reguły dla TypeScriptu
i Reacta) oraz Prettier (jednolite formatowanie); oba narzędzia uruchamia
polecenie \texttt{npm run lint}. Wspólne ustawienia kompilatora TypeScript
(tryb \texttt{strict}) znajdują się w pliku \texttt{tsconfig.base.json},
rozszerzanym przez każdy pakiet.

\section{Opis implementacji}
\label{sec:implementacja}

Kolejne podrozdziały opisują frontend, backend, bazę danych, konteneryzację,
testy i monitoring.

\subsection{Frontend}
\label{sec:frontend}

Frontend to aplikacja jednostronicowa (SPA) napisana w React~19 i TypeScripcie
\cite{react}, budowana narzędziem Vite \cite{vite}. W Dockerze gotowe pliki
serwuje nginx (podrozdział~\ref{sec:docker}). Aplikacja komunikuje się z API
pod tym samym adresem (\texttt{/api}), więc nie wymaga konfiguracji CORS.

\subsubsection{Struktura aplikacji}

\begin{itemize}
  \item \texttt{api/}: klient HTTP (dodaje token, zamienia odpowiedzi błędów na
        wyjątek \texttt{ApiError} z kodem i błędami pól) oraz funkcje wywołujące
        poszczególne endpointy;
  \item \texttt{auth/}: stan zalogowania (kontekst Reacta), zapis sesji
        i strażnicy tras;
  \item \texttt{pages/}: ekrany aplikacji; \texttt{layout/}: układ stron;
        \texttt{components/}: wspólne elementy (pole formularza, komunikat);
  \item \texttt{i18n/texts.ts}: wszystkie teksty interfejsu po polsku w jednym
        module, więc komponenty nie zawierają napisów na stałe.
\end{itemize}

Nawigacją zarządza React Router, a adresy są po polsku (\texttt{/logowanie},
\texttt{/rejestracja}, \texttt{/profil}). Dane z serwera pobiera i przechowuje
w pamięci podręcznej biblioteka TanStack Query. Błędy klienta (4xx) nie są
ponawiane, bo ponowienie dałoby tę samą odpowiedź, a błędy sieci i serwera
są ponawiane jeden raz.

\subsubsection{Logowanie i sesja}

Po zalogowaniu token JWT, dane użytkownika i czas wygaśnięcia tokenu są
zapisywane w \texttt{localStorage}, dzięki czemu sesja przetrwa odświeżenie
strony i ponowne uruchomienie aplikacji zainstalowanej jako PWA. Przy starcie
zapisana sesja jest sprawdzana zapytaniem \texttt{GET /api/auth/me}. Każda
odpowiedź 401 na żądanie z tokenem (token wygasł albo konto usunięto) wylogowuje
użytkownika z komunikatem ,,Sesja wygasła. Zaloguj się ponownie.''. Aplikacja
wylogowuje też sama w chwili wygaśnięcia tokenu. Przechowywanie tokenu
w \texttt{localStorage} oznacza, że skrypt wstrzyknięty na stronę (atak XSS)
mógłby go odczytać. Ryzyko ogranicza React, który domyślnie zabezpiecza
wyświetlane dane, brak zewnętrznych skryptów i krótki czas życia tokenu
(1~godzina). Alternatywą byłoby ciasteczko \texttt{HttpOnly}, które wymaga
jednak dodatkowej ochrony przed atakami CSRF.

Strażnik \texttt{RequireAuth} przekierowuje niezalogowanego użytkownika do
ekranu logowania i zapamiętuje adres, pod który chciał wejść, a po zalogowaniu
kieruje go z powrotem pod ten adres. Strażnik \texttt{GuestOnly} nie pokazuje
ekranów logowania i rejestracji osobie już zalogowanej.

\subsubsection{Formularze}

Formularze obsługuje biblioteka react-hook-form ze schematami Zod
z pakietu \texttt{shared}. Te same reguły i te same polskie komunikaty
obowiązują więc w przeglądarce i w API. Niepoprawny formularz nie jest nawet
wysyłany, a błędy pojawiają się przy polach. Błędy, które zna tylko serwer
(np. zajęty login, odpowiedź 409), są przypisywane do właściwych pól na
podstawie \texttt{details} z odpowiedzi API. Po rejestracji aplikacja
automatycznie loguje użytkownika. Pola formularza są opisane etykietami,
a pole z błędem ma atrybut \texttt{aria-invalid} i powiązany komunikat
(\texttt{aria-describedby}), więc błąd odczytują też czytniki ekranu.

\subsubsection{Układ dla urządzeń mobilnych}

Style są pisane w podejściu \emph{mobile first}: podstawowe reguły dotyczą
telefonów, a szersze ekrany dostają poprawki w media query. Na telefonie
nawigacja główna jest paskiem przyklejonym do dolnej krawędzi ekranu, w zasięgu
kciuka. Wszystkie elementy dotykowe (przyciski, pola, odnośniki nawigacji) mają
co najmniej 48~pikseli wysokości, więcej niż zalecane minimum 44~pikseli.
Układ uwzględnia wycięcia ekranu (\texttt{env(safe-area-inset-*)}). Na ekranie
o szerokości od 768~pikseli nawigacja przechodzi pod górny pasek.

\douzupelnienia{lista i szczegóły dokumentów, profil, obsługa aparatu,
kadrowanie i kompresja (rozmiar przed i po), PWA, zrzuty ekranu z telefonu;
Etapy 8--10}

\subsection{Backend}
\label{sec:backend}

Backend to serwer REST API napisany w TypeScripcie na platformie Node.js
z frameworkiem Express 5 \cite{nodejs,express}. Kod jest kompilowany do modułów
ECMAScript i uruchamiany w kontenerze opisanym w podrozdziale~\ref{sec:docker}.

\subsubsection{Warstwy aplikacji}

Kod podzielono na warstwy, z których każda korzysta tylko z warstwy leżącej
bezpośrednio pod nią:

\begin{itemize}
  \item \textbf{trasy} (\texttt{routes/}) przypisują metody i ścieżki HTTP
        do funkcji kontrolera;
  \item \textbf{kontrolery} (\texttt{controllers/}) odczytują dane z żądania,
        wywołują serwis i budują odpowiedź HTTP (kod statusu, treść JSON);
  \item \textbf{serwisy} (\texttt{services/}) zawierają logikę biznesową
        i nie znają protokołu HTTP;
  \item \textbf{repozytoria} (\texttt{repositories/}) jako jedyne komunikują się
        z bazą danych (Prisma) i magazynem plików (API S3).
\end{itemize}

Aplikację buduje funkcja \texttt{createApp(deps)}, która otrzymuje gotowe
zależności: konfigurację, logger, klienta bazy danych i klienta S3. Tworzy je
dopiero plik \texttt{server.ts}, który następnie uruchamia serwer HTTP. Dzięki
temu testy mogą zbudować aplikację z innymi zależnościami, np. z adresem
nieistniejącej bazy danych, i sprawdzić zachowanie systemu przy awarii
(podrozdział~\ref{sec:testy}).

\subsubsection{Walidacja konfiguracji}

Przy starcie zmienne środowiskowe są sprawdzane schematem biblioteki Zod
\cite{zod}. Sprawdzane są m.in.: format adresu bazy danych, poprawność adresu
URL magazynu MinIO, minimalna długość sekretu JWT (32 znaki) oraz to, czy
klucz główny szyfrowania jest zapisem base64 dokładnie 32 bajtów. Puste
wartości są traktowane jak brakujące. Przy błędnej konfiguracji proces wypisuje
listę \emph{wszystkich} nieprawidłowych zmiennych (bez ich wartości, żeby nie
ujawnić sekretów w logach) i kończy się kodem~1. Kontener z błędną
konfiguracją zatrzymuje się więc od razu i widać przyczynę, zamiast działać
i zgłaszać błędy przy pierwszym żądaniu.

\subsubsection{Obsługa błędów}

Wszystkie błędy trafiają do jednego middleware'u obsługi błędów, który zwraca
je w jednolitym formacie:

\begin{small}
\begin{verbatim}
{ "error": { "code": "VALIDATION_ERROR",
             "message": "Request validation failed",
             "details": [ { "path": "email", "message": "..." } ] } }
\end{verbatim}
\end{small}

Pole \texttt{code} jest stałym identyfikatorem, na podstawie którego frontend
wyświetla komunikat po polsku. Pole \texttt{message} jest opisem dla
programisty, a \texttt{details} zawiera np. błędy poszczególnych pól formularza.
Błędy przewidziane przez aplikację są klasami pochodnymi \texttt{AppError}
(tabela~\ref{tab:bledy}). Każdy inny wyjątek zamieniany jest na odpowiedź
500 \texttt{INTERNAL\_ERROR} z ogólnym komunikatem, a szczegóły (np. treść
błędu bazy danych) trafiają wyłącznie do logu. Express 5 przekazuje do tego
middleware'u również wyjątki z funkcji asynchronicznych, więc kontrolery nie
potrzebują bloków \texttt{try/catch}.

\begin{table}[htbp]
\centering
\caption{Kody błędów API}
\label{tab:bledy}
\small
\begin{tabular}{|l|c|p{0.5\textwidth}|}
\hline
\textbf{Kod} & \textbf{HTTP} & \textbf{Znaczenie} \\
\hline
\texttt{VALIDATION\_ERROR} & 400 & Niepoprawne dane wejściowe (\texttt{ValidationError}) \\
\hline
\texttt{INVALID\_JSON} & 400 & Treść żądania nie jest poprawnym JSON-em \\
\hline
\texttt{FORBIDDEN} & 403 & Brak uprawnień do operacji (\texttt{ForbiddenError}) \\
\hline
\texttt{NOT\_FOUND} & 404 & Nieznana ścieżka lub zasób (\texttt{NotFoundError}) \\
\hline
\texttt{PAYLOAD\_TOO\_LARGE} & 413 & Treść JSON większa niż 1~MB \\
\hline
\texttt{INTERNAL\_ERROR} & 500 & Nieoczekiwany błąd serwera \\
\hline
\end{tabular}
\end{table}

\subsubsection{Logowanie zdarzeń}

Logi są zapisywane biblioteką pino \cite{pino} jako jeden obiekt JSON na
linię. Taki format łatwo przeszukiwać i zbierać centralnie
(podrozdział~\ref{sec:monitoring}). Każde żądanie dostaje identyfikator
(\texttt{requestId}): jeśli klient przesłał poprawny nagłówek
\texttt{X-Request-Id}, jest on używany ponownie, w przeciwnym razie
generowany jest nowy UUID. Identyfikator trafia do nagłówka odpowiedzi i do
każdego wpisu w logu powstałego podczas obsługi żądania. Po zakończeniu
żądania zapisywany jest jeden wpis z metodą, adresem, kodem odpowiedzi i czasem
obsługi, na poziomie \texttt{warn} dla błędów klienta i \texttt{error} dla
błędów serwera. Nagłówki \texttt{Authorization} i \texttt{Cookie} są w logach
zamaskowane. Zapytania kontrolne Dockera do \texttt{/health}, wykonywane co
10~sekund, nie są logowane. W trybie deweloperskim logi są formatowane
czytelnie przez \texttt{pino-pretty}.

\subsubsection{Punkty kontroli stanu}

Backend udostępnia dwa punkty kontroli stanu (tabela~\ref{tab:endpointy}),
dostępne pod ścieżką \texttt{/health} (dla Dockera) oraz \texttt{/api/health}
(przez serwer nginx).

\begin{itemize}
  \item \textbf{Liveness} (\texttt{GET /health}) odpowiada kodem 200, dopóki
        proces obsługuje żądania HTTP. Celowo nie sprawdza bazy danych:
        awaria bazy nie powinna powodować restartu działającego backendu.
  \item \textbf{Readiness} (\texttt{GET /health/ready}) równolegle sprawdza
        PostgreSQL (zapytanie \texttt{SELECT 1}) i MinIO (\texttt{HeadBucket}
        na buckecie dokumentów), każde z limitem 3~sekund. Zwraca 200 i
        \texttt{\{"status":"ready"\}} albo 503 i \texttt{\{"status":"not\_ready"\}},
        a w polu \texttt{checks} status każdej zależności (\texttt{up} lub
        \texttt{down}). Przyczyny awarii są zapisywane w logu, ale nie są
        zwracane klientowi, żeby nie ujawniać wewnętrznych adresów.
\end{itemize}

\begin{table}[htbp]
\centering
\caption{Endpointy API}
\label{tab:endpointy}
\small
\begin{tabular}{|l|l|p{0.4\textwidth}|}
\hline
\textbf{Metoda} & \textbf{Ścieżka} & \textbf{Opis} \\
\hline
GET & \texttt{/health}, \texttt{/api/health} & Liveness -- proces działa \\
\hline
GET & \texttt{/health/ready}, \texttt{/api/health/ready} & Readiness -- baza danych i MinIO dostępne \\
\hline
\end{tabular}
\end{table}

\douzupelnienia{endpointy uwierzytelniania, profilu i dokumentów,
szyfrowanie kopertowe (diagram sekwencji uploadu), wersjonowanie, autoryzacja;
Etapy 4--6}

\subsection{Baza danych}
\label{sec:baza-danych}

Dane aplikacji są przechowywane w bazie PostgreSQL 16 \cite{postgresql},
a dostęp do nich zapewnia Prisma ORM w wersji 7 \cite{prisma}. Schemat bazy
jest zdefiniowany w jednym pliku \path{backend/prisma/schema.prisma}.
Na jego podstawie Prisma generuje klienta z typami TypeScript, więc zapytania
do nieistniejących kolumn i błędy typów wykrywa już kompilator. Prisma 7 nie
korzysta z osobnego silnika zapytań; łączy się z bazą przez sterownik
\texttt{pg} (adapter \texttt{@prisma/adapter-pg}), z limitem 5~sekund na
nawiązanie połączenia.

\subsubsection{Migracje}

Każda zmiana schematu jest zapisywana jako migracja SQL w katalogu
\path{backend/prisma/migrations/} i wersjonowana w repozytorium razem
z kodem. Programista tworzy migrację poleceniem
\verb|npm run db:migrate -w backend -- --name <nazwa>|. W środowisku Docker
migracje stosuje jednorazowa usługa \texttt{migrate}
(\texttt{prisma migrate deploy}), która uruchamia się, gdy baza danych jest
gotowa, a backend startuje dopiero po jej poprawnym zakończeniu. Obraz tej
usługi powstaje z etapu budowania, bo tylko on zawiera narzędzie wiersza
poleceń Prisma. Obraz produkcyjny backendu go nie zawiera.

\subsubsection{Model danych}

Pierwsza migracja (\texttt{init}) tworzy tabelę \texttt{users} z kontami
użytkowników (tabela~\ref{tab:users}). Identyfikatory są typu UUID, więc nie
da się ich odgadnąć, a to ma znaczenie przy adresach dokumentów w API.

\begin{table}[htbp]
\centering
\caption{Tabela \texttt{users}}
\label{tab:users}
\small
\begin{tabular}{|l|l|p{0.45\textwidth}|}
\hline
\textbf{Kolumna} & \textbf{Typ} & \textbf{Opis} \\
\hline
\texttt{id} & \texttt{UUID} & Klucz główny \\
\hline
\texttt{username} & \texttt{VARCHAR(30)} & Login, unikalny \\
\hline
\texttt{email} & \texttt{VARCHAR(254)} & Adres e-mail, unikalny \\
\hline
\texttt{password\_hash} & \texttt{TEXT} & Skrót hasła (argon2), nigdy samo hasło \\
\hline
\texttt{created\_at} & \texttt{TIMESTAMPTZ} & Data utworzenia konta \\
\hline
\texttt{updated\_at} & \texttt{TIMESTAMPTZ} & Data ostatniej zmiany \\
\hline
\end{tabular}
\end{table}

\douzupelnienia{diagram ERD i tabele profilu kierowcy, dokumentów, wersji
i historii zmian; Etapy 5--6}

\subsection{Docker}
\label{sec:docker}

\douzupelnienia{opis docker-compose.yml, wieloetapowe obrazy, sieci i wolumeny; Etap 2}

\subsection{Testy}
\label{sec:testy}

\douzupelnienia{strategia (najpierw testy funkcjonalne), narzędzia, przykładowe scenariusze, raport pokrycia, przykład naprawy błędu metodą test, potem poprawka; Etapy 3--12}

\subsection{Monitoring}
\label{sec:monitoring}

\douzupelnienia{architektura monitoringu, dashboardy Grafany, reguły alertów, watchdogi, wyniki testu chaos-test.sh; Etap 11}


\section{Podsumowanie i wnioski}
\label{sec:podsumowanie}

\subsection{Realizacja celu}

Zrealizowano wszystkie wymagania funkcjonalne z rozdziału~\ref{sec:cel}.
Kierowca zakłada konto i loguje się, a żądania są uwierzytelniane tokenem JWT.
Uzupełnia profil, do którego przypisane są jego dokumenty. Dokument fotografuje
aparatem telefonu, kadruje i obraca zdjęcie, a aplikacja kompresuje je przed
wysłaniem (zdjęcie 12~Mpx o rozmiarze 3,7~MB zajmuje po kompresji około
0,7--0,8~MB). Dokumenty mają typy, statusy z regułami przejść, kolejne wersje
pliku i historię zmian. Można je wyszukiwać i filtrować.

Wymagania niefunkcjonalne zrealizowano następująco:
\begin{itemize}
  \item \textbf{Bezpieczeństwo}: pliki są szyfrowane algorytmem AES-256-GCM
        z osobnym kluczem dla każdego pliku (szyfrowanie kopertowe), hasła są
        przechowywane jako skróty argon2id, a nieudane logowania ograniczane.
        Użytkownik ma dostęp tylko do własnych dokumentów, co sprawdzają testy.
  \item \textbf{Obsługa mobilna}: responsywny interfejs z dolną nawigacją,
        instalowalny jako PWA i uruchamiający się bez sieci.
  \item \textbf{Łatwe wdrożenie}: cały system, łącznie z monitoringiem,
        uruchamia polecenie \texttt{docker compose up}. Migracje bazy
        wykonuje osobny kontener przed startem backendu.
  \item \textbf{Obserwowalność i niezawodność}: metryki Prometheus, cztery
        dashboardy Grafany, 12 reguł alertów, centralne logi w Loki oraz
        automatyczne restarty po awarii i po zawieszeniu procesu, sprawdzone
        testem chaosu.
  \item \textbf{Jakość}: 158 testów backendu (głównie funkcjonalnych, na
        prawdziwej bazie i magazynie plików), 90 testów pakietu wspólnego,
        46 testów jednostkowych frontendu i 31 scenariuszy Playwright
        w emulacji telefonu. Pokrycie kodu przekracza 93~\% instrukcji
        w każdym pakiecie (frontend mierzony testami w przeglądarce).
\end{itemize}

\subsection{Napotkane problemy i wnioski}

Najcenniejszym wnioskiem z projektu jest skuteczność testów funkcjonalnych.
Wykryły one kilka błędów, których nie ujawniłyby testy jednostkowe ani
pobieżne ręczne sprawdzenie:
\begin{itemize}
  \item po odświeżeniu strony pierwsze zapytania wysyłane były bez tokenu,
        bo efekty komponentów potomnych w React wykonują się przed efektem
        dostawcy kontekstu (poprawka: token przechowywany poza drzewem
        komponentów);
  \item strumień aparatu wyłączał się zaraz po uruchomieniu, a użytkownik
        widziałby zamrożoną klatkę (poprawka: efekt uruchamiany raz przy
        wyświetleniu widoku);
  \item po wylogowaniu i zalogowaniu innego kierowcy aplikacja wracała na
        stronę dokumentu poprzedniego użytkownika;
  \item metryki HTTP miały błędne etykiety tras, bo Express przywraca
        ścieżkę bazową routera przed zakończeniem odpowiedzi;
  \item pomiar opóźnienia pętli zdarzeń tracił próbkę po wyzerowaniu
        histogramu, więc zawieszenie procesu nie byłoby wykryte.
\end{itemize}
Każdy z tych błędów najpierw odtworzono testem, który nie przechodził,
a dopiero potem poprawiono kod. Testy pozostają w zestawie i chronią przed
powrotem błędu.

Drugi wniosek dotyczy środowiska testowego. Testy w przeglądarce okazały się
wrażliwe na obciążenie komputera: haszowanie haseł argon2 i przetwarzanie
zdjęć 12~Mpx konkurowały o procesor z przeglądarkami. Stabilność zapewniło
uruchamianie testów na zbudowanej aplikacji (zamiast serwera deweloperskiego)
i kolejne, a nie równoległe, wykonywanie scenariuszy.

Trzeci wniosek dotyczy monitoringu: sama konfiguracja kontroli stanu nie
gwarantuje samonaprawy. Test chaosu pokazał, że Docker nie restartuje
kontenera zatrzymanego poleceniem \texttt{docker kill}, bo traktuje je jak
ręczne zatrzymanie, a zawieszonego, ale działającego procesu nie restartuje
wcale. Dopiero dodatkowy kontener \emph{autoheal} i test z rzeczywistym
zawieszeniem procesu potwierdziły, że system sam wraca do działania.

\subsection{Ograniczenia}

Aparat sprawdzono tylko w emulacji (wirtualna kamera Chromium), bez
prawdziwego telefonu. Procedura testu na urządzeniu z Androidem jest opisana
w pliku README. Potok ciągłej integracji jest przygotowany, ale nie był
uruchomiony, bo repozytorium nie ma zdalnej kopii. Aplikacja ma jedną rolę
użytkownika, więc zmiany statusu, które w praktyce należałyby do biura
(akceptacja, odrzucenie), wykonuje sam kierowca.

\subsection{Możliwości rozwoju}

\begin{itemize}
  \item \textbf{Lepsze zdjęcia dokumentów}: automatyczne wykrywanie krawędzi
        kartki i korekcja perspektywy, jak w aplikacjach skanujących
        (podrozdział~\ref{sec:przeglad}).
  \item \textbf{Rozpoznawanie tekstu (OCR)}: odczyt numeru dokumentu,
        nadawcy i odbiorcy ze zdjęcia i wypełnianie nimi formularza.
  \item \textbf{Panel dyspozytora}: rola pracownika biura, który akceptuje
        lub odrzuca dokumenty kierowców i widzi dokumenty całej floty.
  \item \textbf{e-CMR i eFTI}: import i eksport elektronicznych listów
        przewozowych oraz integracja z platformami eFTI, które od 2027~r.
        mają być w pełni honorowane przez organy kontrolne.
  \item \textbf{Praca bez sieci}: kolejka wysyłek zapisywana na urządzeniu
        w postaci zaszyfrowanej i wysyłana po odzyskaniu połączenia, co ma
        znaczenie na trasach bez zasięgu.
  \item \textbf{Zarządzanie kluczami}: przechowywanie klucza głównego
        w usłudze KMS i jego okresowa rotacja (szyfrowanie kopertowe
        pozwala zmienić klucz główny bez ponownego szyfrowania plików).
\end{itemize}


\spistabel
\spisrysunkow
\printbibliography

\end{document}
